% Options for packages loaded elsewhere
\PassOptionsToPackage{unicode}{hyperref}
\PassOptionsToPackage{hyphens}{url}
\documentclass[
  11pt,
]{article}
\usepackage{xcolor}
\usepackage[margin=1in]{geometry}
\usepackage{amsmath,amssymb}
\setcounter{secnumdepth}{-\maxdimen} % remove section numbering
\usepackage{iftex}
\ifPDFTeX
  \usepackage[T1]{fontenc}
  \usepackage[utf8]{inputenc}
  \usepackage{textcomp} % provide euro and other symbols
\else % if luatex or xetex
  \usepackage{unicode-math} % this also loads fontspec
  \defaultfontfeatures{Scale=MatchLowercase}
  \defaultfontfeatures[\rmfamily]{Ligatures=TeX,Scale=1}
\fi
\usepackage{lmodern}
\ifPDFTeX\else
  % xetex/luatex font selection
\fi
% Use upquote if available, for straight quotes in verbatim environments
\IfFileExists{upquote.sty}{\usepackage{upquote}}{}
\IfFileExists{microtype.sty}{% use microtype if available
  \usepackage[]{microtype}
  \UseMicrotypeSet[protrusion]{basicmath} % disable protrusion for tt fonts
}{}
\makeatletter
\@ifundefined{KOMAClassName}{% if non-KOMA class
  \IfFileExists{parskip.sty}{%
    \usepackage{parskip}
  }{% else
    \setlength{\parindent}{0pt}
    \setlength{\parskip}{6pt plus 2pt minus 1pt}}
}{% if KOMA class
  \KOMAoptions{parskip=half}}
\makeatother
\usepackage{longtable,booktabs,array}
\newcounter{none} % for unnumbered tables
\usepackage{calc} % for calculating minipage widths
% Correct order of tables after \paragraph or \subparagraph
\usepackage{etoolbox}
\makeatletter
\patchcmd\longtable{\par}{\if@noskipsec\mbox{}\fi\par}{}{}
\makeatother
% Allow footnotes in longtable head/foot
\IfFileExists{footnotehyper.sty}{\usepackage{footnotehyper}}{\usepackage{footnote}}
\makesavenoteenv{longtable}
\setlength{\emergencystretch}{3em} % prevent overfull lines
\providecommand{\tightlist}{%
  \setlength{\itemsep}{0pt}\setlength{\parskip}{0pt}}
\usepackage{bookmark}
\IfFileExists{xurl.sty}{\usepackage{xurl}}{} % add URL line breaks if available
\urlstyle{same}
\hypersetup{
  hidelinks,
  pdfcreator={LaTeX via pandoc}}

\author{}
\date{}

\begin{document}

\section{Thought Experiments through the First Three Chapters of
``Foundations of Quantum
Theory''}\label{thought-experiments-through-the-first-three-chapters-of-foundations-of-quantum-theory}

\subsection{------From the Identification Wall in Measurement to the
Algebra of
Observables}\label{from-the-identification-wall-in-measurement-to-the-algebra-of-observables}

\textbf{Author}: Noriaki Kihara, WF System Co., Ltd., ORCID:
\href{https://orcid.org/0009-0004-6753-4020}{0009-0004-6753-4020}
\textbf{Version}: v2.0 \textbf{Date}: May 2026 \textbf{Concept DOI}:
\href{https://doi.org/10.5281/zenodo.20391522}{10.5281/zenodo.20391522}
\textbf{v2 DOI}:
\href{https://doi.org/10.5281/zenodo.20392427}{10.5281/zenodo.20392427}
\textbf{v1 DOI}:
\href{https://doi.org/10.5281/zenodo.20391523}{10.5281/zenodo.20391523}
\textbf{License}: CC BY 4.0

\subsubsection{Abstract}\label{abstract}

Against the background of the framework of complex Hilbert spaces,
observables, and measurement developed in Chapters 1 to 3 of Akira
Shimizu's \emph{New Edition: Foundations of Quantum Theory} {[}1{]},
this paper organizes five stepwise thought experiments to unify the
treatment of measurement precision, the uncertainty relation, wave
packets, quantum correlations, and the algebra of observables. The
starting point is the ``identification wall'' of classical measurement
theory, and the destination is the algebraic picture in which ``physical
quantities satisfying an uncertainty relation are different projections
extracted from the same wave packet.'' The present paper does not
contradict the mathematical system of standard quantum theory; rather,
it offers one reading that supplements its conceptual outlook.

\begin{center}\rule{0.5\linewidth}{0.5pt}\end{center}

\subsection{Thought Experiment I: The Identification Wall in
Measurement}\label{thought-experiment-i-the-identification-wall-in-measurement}

\subsubsection{Setting}\label{setting}

Suppose a physical quantity \emph{A} takes real values. We measure the
distance \emph{L} between two points \emph{A}₁ and \emph{A}₂ using a
measure based on the concept of distance, which returns only discrete
values with a minimum spacing of Δ.

At each measurement, one of the two adjacent lattice points \{\emph{k}Δ,
(\emph{k}+1)Δ\} bracketing the true value is returned with probability
depending on the position \emph{r} of the true value within the cell:

P(up) = \emph{r}/Δ, P(down) = 1 − \emph{r}/Δ

(position-dependent dithering).

\subsubsection{Behavior of the N-fold
Average}\label{behavior-of-the-n-fold-average}

When \emph{A}₁ = \emph{A}₂, the single-shot distance \emph{D} = ε₂ − ε₁
takes values \{−Δ, 0, +Δ\} with probabilities \{1/4, 1/2, 1/4\}.

\begin{itemize}
\tightlist
\item
  E{[}\emph{D}{]} = 0
\item
  Var{[}\emph{D}{]} = Δ²/2
\item
  Standard deviation of the N-fold average: σ\_N = Δ/√(2N)
\end{itemize}

As N → ∞, σ\_N → 0, and L̂\_N → \emph{L} with probability 1.

\subsubsection{Condition for
Distinguishability}\label{condition-for-distinguishability}

For \emph{A}₁ and \emph{A}₂ to be recognized as two distinct objects,
the two must belong to different Δ cells. If they fall within the same
cell, the measure cannot distinguish ``one point or two,'' and the
quantity ``distance between two points'' cannot be constituted.

Therefore the condition for the problem of measuring \emph{L} to be
well-posed is

\emph{L} ≥ Δ.

The region \emph{L} \textless{} Δ is not ``buried in measurement
error''; it is a region in which the very concept of ``two points'' is
not defined.

\subsubsection{Conclusion}\label{conclusion}

\begin{quote}
When a measure based on distance is used to measure distance,
measurement at scales below the measure's spacing Δ is
\textbf{undefined} prior to any invocation of the uncertainty principle.
Δ functions not as a resolution but as the threshold of individuation.

Within the range \emph{L} ≥ Δ, arbitrary precision is attainable by
repeated measurement regardless of the absolute value of Δ.
\end{quote}

The present argument is related to but independent of the discussions of
minimum length at the Planck scale in quantum gravity {[}10{]}; it holds
in general for any measure based on the concept of distance.

\begin{center}\rule{0.5\linewidth}{0.5pt}\end{center}

\subsection{Thought Experiment II: The Locus of
Fluctuation}\label{thought-experiment-ii-the-locus-of-fluctuation}

\subsubsection{Setting}\label{setting-1}

Consider the reverse situation. The measure has infinite precision and
returns its input value unchanged. The physical quantity \emph{A}
itself, however, fluctuates about its true value \emph{A'} according to

\emph{A}\_i = \emph{A'} + ξ\_i, ξ\_i \textasciitilde{} Uniform{[}−Δ/2,
+Δ/2{]}.

We assume \emph{L} \textgreater{} Δ.

\subsubsection{N-fold Average}\label{n-fold-average}

\begin{itemize}
\tightlist
\item
  E{[}L̂\_N{]} = \emph{A'}
\item
  Var{[}\emph{L}\_i{]} = Δ²/12
\item
  σ\_N = Δ/√(12N) → 0
\end{itemize}

\subsubsection{Comparison with Thought Experiment
I}\label{comparison-with-thought-experiment-i}

Instantiating both cases with true value 7.7 and spacing Δ = 1, each
yields the same empirical distribution: out of 100 measurements, roughly
30 return 7 and 70 return 8, with the average converging to 7.7.

{\def\LTcaptype{none} % do not increment counter
\begin{longtable}[]{@{}
  >{\raggedright\arraybackslash}p{(\linewidth - 4\tabcolsep) * \real{0.3333}}
  >{\raggedright\arraybackslash}p{(\linewidth - 4\tabcolsep) * \real{0.3333}}
  >{\raggedright\arraybackslash}p{(\linewidth - 4\tabcolsep) * \real{0.3333}}@{}}
\toprule\noalign{}
\begin{minipage}[b]{\linewidth}\raggedright
\end{minipage} & \begin{minipage}[b]{\linewidth}\raggedright
Fluctuation on the measure side
\end{minipage} & \begin{minipage}[b]{\linewidth}\raggedright
Fluctuation on the quantity side
\end{minipage} \\
\midrule\noalign{}
\endhead
\bottomrule\noalign{}
\endlastfoot
Single-shot output space & Discrete \{\emph{k}Δ, (\emph{k}+1)Δ\} &
Continuous \\
Probability distribution & Position-dependent Bernoulli & Additive
uniform \\
N → ∞ convergence target & True value & True value \\
Distinguishability from observed data & Impossible & \\
\end{longtable}
}

\subsubsection{Two-sided Fluctuation}\label{two-sided-fluctuation}

When fluctuations ξ (width δ) on the quantity side and η (width σ) on
the measure side coexist independently:

\begin{itemize}
\tightlist
\item
  Observed value: \emph{L}\_i = \emph{A'} + ξ\_i + η\_i
\item
  σ\_N = √(Var{[}ξ{]} + Var{[}η{]}) / √N → 0
\end{itemize}

The 1/√N decay persists, but separating δ and σ from a single observed
series requires additional assumptions external to the observation
(independent calibration of the measure, time-series decomposition,
prior knowledge of distribution shapes, etc.).

\subsubsection{Conclusion}\label{conclusion-1}

\begin{quote}
Defining the observed fluctuation magnitude as Δ\_obs, one cannot
determine from observed data alone whether it originates from the
quantization width of the measure, from the fluctuation of the quantity
itself, or from a combination of both.

Rather than ``convergence to the objective true value,'' the precise
statement is ``convergence to the effective true value defined by the
observer,'' and the 1/√N decay applies rigorously in this sense.
\end{quote}

\begin{center}\rule{0.5\linewidth}{0.5pt}\end{center}

\subsection{Thought Experiment III: Wavelength Representation of
Momentum and the Uncertainty
Relation}\label{thought-experiment-iii-wavelength-representation-of-momentum-and-the-uncertainty-relation}

\subsubsection{Setting}\label{setting-2}

We extend the object of Thought Experiment II to a quantum-theoretic
wave packet. For position \emph{x}, consider the center value \emph{x}₁
and the spread δ\emph{x} of the wave packet. Here ``center value'' is
merely a label of the packet distribution; we do not take the position
that there exists an objective true value external to observation.

Similarly for momentum \emph{p}, consider the center value \emph{p}₁ and
the spread δ\emph{p}.

\subsubsection{Rewriting via Wavelength}\label{rewriting-via-wavelength}

From the de Broglie relation {[}2{]}:

\[p = \frac{h}{\lambda} = \hbar k, \quad k = \frac{2\pi}{\lambda}\]

the momentum fluctuation is

\[\delta p = \frac{h}{\lambda^2}\delta\lambda = \hbar \, \delta k.\]

Rewriting the uncertainty relation Δ\emph{x}Δ\emph{p} ≥ ℏ/2 {[}7{]} in
terms of the wavenumber:

\[\Delta x \cdot \Delta k \geq \frac{1}{2}.\]

\subsubsection{Observation}\label{observation}

In this form, Planck's constant ℏ disappears from both sides. What
remains is a mathematical inequality of Fourier transforms concerning
the product of position and wavenumber.

Planck's constant ℏ functions as a dimensional conversion factor
connecting the unit of position (m) and the unit of momentum (kg·m/s). A
similar view that treats Planck's constant not as an observable quantity
but as a convention of human unit-choice has been proposed by Ralston
{[}8{]}.

Qualitatively:

\begin{itemize}
\tightlist
\item
  Pinning position precisely ⟺ narrowing the position wave packet ⟺
  broadening the wavelength distribution ⟺ making momentum vague
\item
  Pinning momentum precisely ⟺ narrowing the wavelength to one ⟺
  spreading over all of space ⟺ making position vague
\end{itemize}

{\def\LTcaptype{none} % do not increment counter
\begin{longtable}[]{@{}lll@{}}
\toprule\noalign{}
Packet form & Spatial spread & Wavelength distribution \\
\midrule\noalign{}
\endhead
\bottomrule\noalign{}
\endlastfoot
Delta function & δ\emph{x} = 0 & All wavelengths \\
Rectangular pulse & Finite & sinc function \\
Pure sine wave & All of space & Single wavelength \\
\end{longtable}
}

\subsubsection{Conclusion}\label{conclusion-2}

\begin{quote}
In the wavenumber representation, the position--momentum uncertainty
relation becomes Δ\emph{x}Δ\emph{k} ≥ 1/2, described as a lower bound on
the product of conjugate quantities in Fourier transforms. Planck's
constant appears as a dimensional conversion factor.
\end{quote}

\begin{center}\rule{0.5\linewidth}{0.5pt}\end{center}

\subsection{Thought Experiment IV: Quantum Correlations as a Composite
Wave
Packet}\label{thought-experiment-iv-quantum-correlations-as-a-composite-wave-packet}

\subsubsection{Setting}\label{setting-3}

Consider a state localized in two spatial regions \emph{A} and \emph{B}
with correlations mediated by a shared conserved quantity (momentum,
energy, etc.). In standard descriptions, this system is treated as a
``two-particle entangled state'' {[}3, 4, 5{]}.

We attempt to describe this instead as ``a state in which a single
composite wave packet is localized in two regions in space.''

\subsubsection{String Analogy}\label{string-analogy}

When a long, taut string undergoing baseline vibration has one end
suddenly fixed, a soliton-like deformation appears at the opposite end.
This is a result of the entire string obeying conservation laws as a
single system; no information has propagated from one end to the other.

\subsubsection{Geometry of the Composite Wave
Packet}\label{geometry-of-the-composite-wave-packet}

\begin{itemize}
\tightlist
\item
  The composite wave packet occupies one area element in phase space
\item
  In configuration space it has localized peaks in two regions \emph{A}
  and \emph{B}
\item
  A local interaction occurs on the \emph{A} side, contracting
  Δ\emph{x}\_A
\item
  Due to area conservation (the lower bound of the Robertson inequality
  {[}6{]}), the conjugate Δ\emph{p}\_A broadens
\item
  To satisfy the conserved quantities of the entire composite packet,
  the spread of the wave packet on the \emph{B} side also undergoes
  geometric deformation
\end{itemize}

In this description, no information transfer from \emph{A} to \emph{B}
is introduced. The observed correlation is a geometric consequence of a
local update to one wave packet being reflected throughout via the
conservation laws.

\subsubsection{The True Nature of the Conserved
Quantity}\label{the-true-nature-of-the-conserved-quantity}

What is conserved in the composite system is not individual wavelengths
or positions, but the product of the spreads of conjugate physical
quantities (an area element in phase space):

\[\Delta x \cdot \Delta p \geq \frac{\hbar}{2}, \quad \Delta A \cdot \Delta B \geq \frac{1}{2}\left|\langle [A,B] \rangle\right|.\]

The left side has the dimension of area; the inequality means a lower
bound on the area.

\subsubsection{Conclusion}\label{conclusion-3}

\begin{quote}
Quantum correlations are describable as ``a state in which one composite
wave packet is localized across multiple regions in space.'' The
observed correlations are geometric consequences of conservation laws
acting on the entire packet, and this description does not require
information transfer between spatially separated systems.
\end{quote}

\begin{center}\rule{0.5\linewidth}{0.5pt}\end{center}

\subsection{Thought Experiment V: The Algebra of
Observables}\label{thought-experiment-v-the-algebra-of-observables}

\subsubsection{Setting}\label{setting-4}

We list several pairs of physical quantities that satisfy the
uncertainty relation Δ\emph{A}Δ\emph{B} ≥
\textbar⟨{[}\emph{A},\emph{B}{]}⟩\textbar/2:

\begin{itemize}
\tightlist
\item
  Position \emph{x} ↔ momentum \emph{p}
\item
  Spin \emph{S}\_x ↔ \emph{S}\_y ↔ \emph{S}\_z
\item
  Polarization (H/V ↔ D/A ↔ R/L)
\item
  Angular momentum \emph{L}\_x ↔ \emph{L}\_y ↔ \emph{L}\_z
\item
  Energy ↔ time
\end{itemize}

All of these share a common structure.

\subsubsection{Common Structure}\label{common-structure}

{\def\LTcaptype{none} % do not increment counter
\begin{longtable}[]{@{}
  >{\raggedright\arraybackslash}p{(\linewidth - 4\tabcolsep) * \real{0.3333}}
  >{\raggedright\arraybackslash}p{(\linewidth - 4\tabcolsep) * \real{0.3333}}
  >{\raggedright\arraybackslash}p{(\linewidth - 4\tabcolsep) * \real{0.3333}}@{}}
\toprule\noalign{}
\begin{minipage}[b]{\linewidth}\raggedright
Pair of quantities
\end{minipage} & \begin{minipage}[b]{\linewidth}\raggedright
Shared structure
\end{minipage} & \begin{minipage}[b]{\linewidth}\raggedright
Role of distinct operators
\end{minipage} \\
\midrule\noalign{}
\endhead
\bottomrule\noalign{}
\endlastfoot
\emph{x} ↔ \emph{p} & Area element in phase space & Direct coordinate ↔
Fourier coordinate \\
\emph{S}\_x ↔ \emph{S}\_y ↔ \emph{S}\_z & Area element in spin phase
space & Projection onto distinct axes \\
Polarization basis & Area element of polarization plane & Projection
onto distinct bases \\
\emph{L}\_x ↔ \emph{L}\_y ↔ \emph{L}\_z & Angular-momentum phase space &
Projection onto distinct axes \\
\emph{E} ↔ \emph{t} & Energy--time phase space & Distinct
representations \\
\end{longtable}
}

\subsubsection{Unified Picture}\label{unified-picture}

In a complex Hilbert space:

{\def\LTcaptype{none} % do not increment counter
\begin{longtable}[]{@{}
  >{\raggedright\arraybackslash}p{(\linewidth - 2\tabcolsep) * \real{0.5000}}
  >{\raggedright\arraybackslash}p{(\linewidth - 2\tabcolsep) * \real{0.5000}}@{}}
\toprule\noalign{}
\begin{minipage}[b]{\linewidth}\raggedright
Formal object
\end{minipage} & \begin{minipage}[b]{\linewidth}\raggedright
Geometric interpretation
\end{minipage} \\
\midrule\noalign{}
\endhead
\bottomrule\noalign{}
\endlastfoot
Eigenvector \textbar ψ⟩ & Wave packet (one area element in phase
space) \\
Operator  & Projection extracting a real value from the area element \\
Eigenvalue \emph{a} & Result of the projection \\
Non-commutativity {[}Â, B̂{]} ≠ 0 & Two projections viewing the same area
element from different directions \\
\end{longtable}
}

\subsubsection{Conclusion}\label{conclusion-4}

\begin{quote}
Physical quantities satisfying an uncertainty relation are described as
distinct projections of the same wave packet (eigenvector). The product
of non-commuting projections is bounded below by the area element of the
wave packet (Robertson inequality).

A physical quantity is a projection of a wave packet. The wave packet
exists as an area element in phase space, and the area is conserved as
the product of projections.
\end{quote}

The area-element picture of the present paper is consistent with de
Gosson's \emph{quantum blobs} {[}9{]} --- the smallest phase-space units
saturating the Robertson--Schrödinger inequality, measured by symplectic
capacity (which has the dimension of area). The present paper can be
regarded as a reading of Chapter 1 of Shimizu's textbook that
reconstructs this mathematically heavy symplectic-geometric framework in
the intuitive form of thought experiments. Related geometrical
formulations include the Kähler-geometric formulation on projective
Hilbert space {[}11{]} and the algebraic approach in which the algebra
of observables is taken as primary {[}12{]}.

\begin{center}\rule{0.5\linewidth}{0.5pt}\end{center}

\subsection{Summary}\label{summary}

Through the five thought experiments, the following structure emerged:

\begin{enumerate}
\def\labelenumi{\arabic{enumi}.}
\tightlist
\item
  \textbf{The identification wall} (I) --- Δ is not a resolution but a
  threshold of individuation. \emph{L} ≥ Δ is a prerequisite for
  distance measurement.
\item
  \textbf{Indistinguishability of the locus of fluctuation} (II) --- Δ
  on the measure side and Δ on the quantity side cannot be distinguished
  from observed data.
\item
  \textbf{Uncertainty relation in the wavenumber representation} (III)
  --- Δ\emph{x}Δ\emph{k} ≥ 1/2 is a mathematical fact of Fourier
  transforms. ℏ appears as a unit conversion factor.
\item
  \textbf{Quantum correlations as a composite wave packet} (IV) ---
  Described as the spatial division of a single wave packet, the
  correlations between spatially separated systems can be explained
  without assuming information transfer.
\item
  \textbf{The algebra of observables} (V) --- Physical quantities
  satisfying an uncertainty relation are distinct projections of the
  same wave packet. The area element in phase space is the conserved
  quantity.
\end{enumerate}

These are consistent with the framework of complex Hilbert spaces,
observables, and measurement introduced in Chapters 1 to 3 of Shimizu's
\emph{New Edition: Foundations of Quantum Theory}, and present one
reading that provides a conceptual outlook on quantum theory. The
present paper does not modify the mathematical predictions of standard
quantum theory; it confines itself to offering a geometric
interpretation of each concept.

\begin{center}\rule{0.5\linewidth}{0.5pt}\end{center}

\subsection{Acknowledgment}\label{acknowledgment}

The present thought experiments were progressively organized through
dialogue with AI. Verbatim records including the trial-and-error process
are publicly available in the ai-chat-logs-open repository
(\texttt{新版量子論の基礎/思考実験(6)*}, \texttt{思考実験(7)*}).

\subsection{References}\label{references}

{[}1{]} Shimizu, A. \emph{New Edition: Foundations of Quantum Theory ---
For an Easy Understanding of Its Essence} (in Japanese), Saiensu-sha,
2003.

{[}2{]} de Broglie, L. (1924) \emph{Recherches sur la théorie des
Quanta}, doctoral thesis, Faculté des Sciences de Paris.

{[}3{]} Einstein, A., Podolsky, B., Rosen, N. (1935) ``Can
Quantum-Mechanical Description of Physical Reality Be Considered
Complete?'' \emph{Physical Review} \textbf{47}, 777--780.

{[}4{]} Schrödinger, E. (1935) ``Discussion of Probability Relations
Between Separated Systems,'' \emph{Mathematical Proceedings of the
Cambridge Philosophical Society} \textbf{31}, 555--563.

{[}5{]} Bell, J. S. (1964) ``On the Einstein-Podolsky-Rosen Paradox,''
\emph{Physics Physique Физика} \textbf{1}, 195--200.

{[}6{]} Robertson, H. P. (1929) ``The Uncertainty Principle,''
\emph{Physical Review} \textbf{34}, 163--164.

{[}7{]} Heisenberg, W. (1927) ``Über den anschaulichen Inhalt der
quantentheoretischen Kinematik und Mechanik,'' \emph{Zeitschrift für
Physik} \textbf{43}, 172--198.

{[}8{]} Ralston, J. P. (2020) ``Quantum Theory without Planck's
Constant,'' \emph{International Journal of Quantum Foundations}
\textbf{6}(3), 1--43. Originally arXiv:1203.5557 (2012).

{[}9{]} de Gosson, M. A. (2013) ``Quantum Blobs,'' \emph{Foundations of
Physics} \textbf{43}, 440--457.

{[}10{]} Garay, L. J. (1995) ``Quantum Gravity and Minimum Length,''
\emph{International Journal of Modern Physics A} \textbf{10}, 145--166.
arXiv:gr-qc/9403008.

{[}11{]} Ashtekar, A. and Schilling, T. A. (1999) ``Geometrical
Formulation of Quantum Mechanics,'' in \emph{On Einstein's Path: Essays
in Honor of Engelbert Schücking} (ed.~A. Harvey), Springer, pp.~23--65.
arXiv:gr-qc/9706069.

{[}12{]} Slavnov, D. A. (2007) ``Algebraic Quantum Theory,''
\emph{Physics of Particles and Nuclei} \textbf{38}, 147--176.

\end{document}
